\chapter{相关工作}

\section{系统组件及其功能概括}

本系统由文档解析、Embedding/Fusion、检索重排、GRPO 决策模型、多模态生成和前端溯源预览组成，共同完成“获取证据、选择模型、生成答案、回到原文档验证”的流程。

\begin{table}[htbp]
  \centering
  \caption{相关技术组件与项目功能对应}
  \label{tab:component_function}
  \footnotesize
  \setlength{\tabcolsep}{4pt}
  \renewcommand{\arraystretch}{1.12}
  \begin{tabularx}{\textwidth}{>{\raggedright\arraybackslash}p{0.22\textwidth}>{\raggedright\arraybackslash}X}
    \toprule
    技术组件 & 功能与作用 \\
    \midrule
    文档解析与 OCR & 将 PDF/图片转成页面、文本、坐标和局部区域，支撑 chunk 构建、PDF 预览和证据高亮。\\
    Embedding 与 Fusion & 将文本、OCR、表格和图像特征转成统一 evidence，使不同模态可以共同检索和生成。\\
    向量检索与 Reranker & 先召回候选证据，再过滤弱相关 chunk，提高送入生成模型的证据质量。\\
    GRPO 决策模型 & 通过强化学习训练决策模型，根据问题和证据选择更合适的 LLM/VLM 与上下文组织方式。\\
    VLM 多模态生成 & 根据重排后的文本与视觉 evidence 构建上下文，并调用选定模型生成答案。\\
    前端溯源预览 & 展示答案对应页面和 bbox 区域，使问答结果可检查、可追溯。\\
    \bottomrule
  \end{tabularx}
\end{table}

\section{GRPO 决策模型}

GRPO 决策模块对应系统架构图中的强化学习微调部分，其核心目标是训练一个 Decide Model，使其能够为不同问题选择更合适的 LLM/VLM 进行回答。该模块不是简单的关键词规则判断，而是通过“模型回答-裁判打分-策略更新”的闭环，让决策模型逐渐学习不同问题与不同模型能力之间的对应关系。

具体而言，Decide Model 先根据用户问题和检索证据选择若干候选模型；候选模型分别生成回答；Judge Model（如 Qwen-Max）根据答案正确性、证据一致性和回答质量给出 reward，例如 0、0.5 或 1；最后使用 GRPO 根据这些奖励更新 Decide Model。这样，系统可以在训练过程中让决策模型学会：什么类型的问题更适合文本 LLM，什么类型的问题更适合多模态 VLM，什么情况下需要依赖局部视觉区域或图表信息。

\section{多模态向量化}

在数据预处理阶段，系统通过版面分析将 PDF 或图片切分为页面、文本块、表格、图像区域和局部视觉 evidence，并根据模态进行向量化。

\begin{itemize}
  \item \textbf{文本}：使用 BGE-M3 对普通段落、OCR 文本和代码片段向量化。
  \item \textbf{表格}：将表格转为 Markdown，保留行列结构后进入索引。
  \item \textbf{图片}：使用轻量视觉特征或 ViT 图像编码器提取图像向量，并保留页码与 bbox。
  \item \textbf{局部区域}：对图表、公式、logo、广告和票据区域保留 crop、OCR anchor 和来源页面。
\end{itemize}

多模态向量化使文本、表格和图像证据可以在统一接口中被召回、排序和使用。

\section{在线检索与决策}

\begin{itemize}
  \item \textbf{混合检索 (Hybrid Retrieval)}：同时搜索文本 chunk、表格 Markdown、OCR 片段和视觉区域描述，获取 Top-K 候选 evidence。
  \item \textbf{Rerank 重排序}：使用 BGE-Reranker-v2-m3 对问题和候选证据二次打分，过滤无关材料。
  \item \textbf{模型决策}：GRPO 决策模型结合问题和候选 evidence 选择回答模型；候选模型生成答案后，由 Judge Model 评分，评分结果继续反馈给 Decide Model。
\end{itemize}

\section{Generation 生成阶段}

\begin{itemize}
  \item \textbf{上下文构建}：将文本、表格 Markdown、OCR 片段、页面图和局部 crop 组织为带编号 evidence。
  \item \textbf{VLM 推理}：对于图表、表格或视觉区域问题，调用 Qwen3-VL-8B-Instruct 进行多模态推理；普通文本问题使用文本 RAG 生成。
  \item \textbf{答案归一与溯源}：对输出进行短答案抽取、数值容差和实体归一，并将证据编号映射回 page、chunk 和 bbox。
\end{itemize}

\section{前端技术栈与系统功能}

前端使用 React、Vite、TailwindCSS 和 Framer Motion 构建，后端通过 FastAPI 提供文件上传、解析进度、chunk 获取、页面图读取、检索问答和文件删除接口。

\begin{itemize}
  \item \textbf{文档管理}：支持上传 PDF/图片，展示解析状态、文件列表和知识库 chunk。
  \item \textbf{问答交互}：根据用户问题自动完成检索、模型决策、重排和生成。
  \item \textbf{连续 PDF 预览}：右侧 viewer 支持上下翻页查看完整文档。
  \item \textbf{证据高亮}：点击 chunk 或答案引用后，根据 bbox 定位并高亮原文档区域。
\end{itemize}

该交互参考 RAGFlow 等文档知识库系统\cite{ragflow}，将答案生成与原文档证据检查放在同一界面中。
